\documentclass[12pt,letterpaper]{beamer}
\usetheme{Copenhagen}
\usecolortheme{seahorse}
\setbeamertemplate{section in toc}{\inserttocsection}

\usepackage[utf8]{inputenc}
\usepackage{amsmath}
\usepackage{amsfonts}
\usepackage{amssymb}
\usepackage{graphicx}
\graphicspath{ {./images/} }
\usepackage{multirow}
\usepackage{hyperref}
\hypersetup{
    colorlinks=true,
    linkcolor=blue,
    filecolor=magenta,      
    urlcolor=cyan,
    pdftitle={Overleaf Example},
    pdfpagemode=FullScreen,
}
\usepackage{minted}

\title[Robotics I]
{ENGR 3421: ROBOTICS I}
\subtitle{Python Advanced}

\author[Zhang, Lin]
{Dr. Lin Zhang}
\institute[UCA] % (optional)
{
  Department of Physics and Astronomy\\
  University of Central Arkansas
}
\date[Robotics1 2021] % (optional)
{September 23, 2021}
\logo{\includegraphics[height=1cm]{../images/uca_bear_logo.png}}


%End of title page configuration block
%------------------------------------------------------------

%------------------------------------------------------------
%The next block of commands puts the table of contents at the beginning of each section and highlights the current section:

\AtBeginSection[]
{
  \begin{frame}
    \frametitle{Outline}
    \tableofcontents[currentsection]
  \end{frame}
}
%------------------------------------------------------------

\begin{document}

%The next statement creates the title page.
\frame{\titlepage}

%---------------------------------------------------------
%This block of code is for the table of contents after the title page
\begin{frame}
\frametitle{Outline}
\tableofcontents
\end{frame}
%---------------------------------------------------------

\section{Class}

\begin{frame}{Object-Oriented Programming}

Object Oriented programming (OOP) is a programming paradigm that relies on the concept of classes and objects. 
It is used to structure a software program into simple, reusable pieces of code blueprints (usually called classes), which are used to create individual instances of objects.
\end{frame}

\begin{frame}[fragile]
    \frametitle{Function w/o Return}
        {\scriptsize
        \begin{minted}{python}
            def forward(motor1, motor2, speed=1):
                """
                Args:
                    motor1: object instantiate from Motor class
                    motor2: object instantiate from Motor class
                    speed: scalar in range [0,1]
                Return:
                    None
                """
                motor1.set_speed(speed)
                motor2.set_speed(speed)
        
        \end{minted}
        }
\end{frame}

\begin{frame}[fragile]
    \frametitle{Function w/ Return}

        {\scriptsize
        \begin{minted}{python}
            def compute_center(ul_coord, ur_coord, lr_coord, ll_coord):
                """
                Args:
                    ul_coord: array with shape (2,) or list with length 2
                    ur_coord: [x, y]
                    lr_coord: e.g. array(321, 456)
                    ll_coord: e.g. [321, 456]
                Return:
                    center_coord: coordinate of center of the box represented 
                    by a list with length of 2.
                """
                mean_x = (ul_coord[0] + \
                          ur_coord[0] + \
                          lr_coord[0] + \
                          ll_coord[0]) / 4
                mean_y = (ul_coord[1] + \
                          ur_coord[1] + \
                          lr_coord[1] + \
                          ll_coord[1]) / 4
                center_coord = [mean_x, mean_y]
        
                return center_coord
        \end{minted}
        }

\end{frame}


\end{document}


